\clearpage\begingroup\SourceLeaf{419b}\ArithmeticHead{88}{L}{STÄCKEL, GAUSS ALS GEOMETER.}
\renewcommand{\thefootnote}{\arabic{footnote})}\setcounter{footnote}{0}
\makeatletter\tagsleft@true\let\veqno\leqno\makeatother
\begin{center}
\textbf{29.}\par\smallskip
\textls[100]{Die Kopenhagener Preisschrift} (1822).
\end{center}
\noindent\hspace*{\parindent}So lebhaft der Anteil war, den \textsc{Gauss} an den Gradmessungen nahm, so
warm er die Nachricht von \textsc{Schumachers} Unternehmen begrüsst hatte, so hat
er sich doch über den Vorschlag, den Meridian durch Hannover fortzusetzen,
zurückhaltend geäussert. »In diesem Augenblick«, schreibt er am 5. Juli 1816,
»kann ich zwar solchen Wunsch in Hannover noch nicht in Anregung bringen,
da erst die Astronomie selbst noch so grosser Unterstützung bedarf: allein ich
bin überzeugt, dass demnächst unsere Regierung, die auch die Wissenschaften
gern unterstützt, dem glorreichen Beispiel Ihres trefflichen Königs folgen
werde« (W. IX, S. 345). In der Tat näherte sich zu dieser Zeit der lange
\clearpage\SourceLeaf{420}\ArithmeticHead{89}{R}{DIE ALLGEMEINE LEHRE VON DEN KRUMMEN FLÄCHEN.}\setcounter{footnote}{0}
\noindent hingezogene Neubau der Göttinger Sternwarte der Vollendung, und \textsc{Gauss} war
im April und Mai 1816 in München gewesen, um mit \textsc{Reichenbach} und \textsc{Steinheil}
wegen der neu zu beschaffenden Messwerkzeuge zu verhandeln. Im
Herbst des Jahres hat er dann seinen Einzug in die Räume gehalten, die er
fast 40 Jahre innehaben sollte.
\par
\textsc{Gaede} hat auf Grund der Akten dargelegt, wie der »welterfahrene und
geschäftsgewandte« \textsc{Schumacher} (Br. G.-Sch. I, S. 190) in jahrelangen Verhandlungen
die Schwierigkeiten überwand, die sich seinem zum »Sollizitieren«
wenig geneigten und geeigneten Freunde (Br. G.-Sch. I, S. 142) entgegenstellten,
bis dieser endlich durch die Kabinettsordre Georgs IV., Königs von
England und Hannover, vom 9. Mai 1820 den Auftrag zur Ausführung der
Gradmessung erhielt\footnote{Vgl. \textsc{Gaede}, \textit{Beiträge zur Kenntnis von Gauss’ praktisch-geodätischen Arbeiten}, Zeitschrift für Vermessungswesen, Bd. 14, 1885; auch als besondere Schrift, Karlsruhe 1885 erschienen.}. Die Messungen im Felde haben fünf Arbeitsjahre,
1821 bis 1825, erfordert, und im Frühjahr 1827 folgte noch die astronomische
Bestimmung des Breitenunterschiedes der Sternwarten zu Göttingen und zu
Altona. Nunmehr wurde durch die Kabinettsordre vom 25. März 1828 die
Triangulation des ganzen Königreichs Hannover befohlen, und \textsc{Gauss} am
14. April vom Ministerium mit der Leitung beauftragt. Wenn er auch an
den Aufnahmen im Feld nicht mehr teilnahm, so erwuchs ihm doch aus den
Messungsergebnissen eine grosse und öde Rechenarbeit, die erst mit dem
Jahre 1848 zum Abschluss gekommen ist. Wiederholt hat \textsc{Gauss} beklagt,
wie sehr er dadurch in seinen wissenschaftlichen Untersuchungen gehemmt
werde. »Gewiss ist, dass wenn meine Lage immer die nämliche bleibt, ich
den grössern Teil meiner früheren theoretischen Arbeiten, denen noch, der
einen mehr, der andern weniger an der Vollendung fehlt, und die von solcher
Art sind, dass Vollendung sich nicht erzwingen lässt, wenn man eben will,
mit ins Grab nehmen werde. Denn etwas Unvollendetes kann und mag ich
einmal nicht geben« (Brief an \textsc{Bessel} vom 15. November 1822, Br. G.-\textsc{Bessel}
S. 410). Nur wer sich in eine solche Lage und Stimmung zu versetzen vermag,
wird verstehen, wie es gekommen ist, dass \textsc{Gauss} nur einen Teil seiner
umfangreichen Untersuchungen über die allgemeine Lehre von den krummen
Flächen ausgearbeitet und bekanntgegeben hat.
\par
Wir verdanken es wiederum \textsc{Schumacher}, dass \textsc{Gauss} mit der Veröffentlichung
\par\smallskip\noindent{\small X\textsubscript{2} Abh. 4.\hfill 12}\par
\clearpage\SourceLeaf{421}\ArithmeticHead{90}{L}{STÄCKEL, GAUSS ALS GEOMETER.}\setcounter{footnote}{0}
\noindent seiner Entdeckungen einen Anfang machte. Wie schon erwähnt
wurde, hatte \textsc{Gauss} in dem Briefe vom 5. Juli 1816 von einer Preisfrage erzählt,
die er für die neue astronomische Zeitschrift vorgeschlagen hatte, die
aber nicht gewählt worden war. »Mir war eine interessante Aufgabe eingefallen«,
schreibt er, »nämlich: \textls[100]{allgemein} eine gegebene Fläche so auf einer
andern (gegebenen) zu projizieren (abzubilden), dass das Bild dem Original in
den kleinsten Teilen ähnlich werde. Ein spezieller Fall ist, wenn die erste
Fläche eine Kugel, die zweite eine Ebene ist. Hier sind die stereographische
und die merkatorische Projektionen partikuläre Auflösungen. Man will aber
die \textls[100]{allgemeine} Auflösung, worunter alle partikulären begriffen sind, für
jede Arten von Flächen. Es soll darüber in dem Journal philomathique bereits
von \textsc{Monge} und \textsc{Poinsot} gearbeitet sein (wie \textsc{Burckhardt} an \textsc{Lindenau}
geschrieben hat), allein da ich nicht genau weiss wo, so habe ich noch nicht
nachsuchen können und weiss daher nicht, ob jener Herren Auflösungen ganz
meiner Idee entsprechen und die Sache erschöpfen« (W. VIII, S. 370)\footnote{Weder in dem Bulletin de la société philomathique noch in den sonstigen Veröffentlichungen von \textsc{Monge} und \textsc{Poinsot} hat sich eine auf die konforme Abbildung bezügliche Stelle finden lassen. Vielleicht hat \textsc{Burckhardt} an \textsc{Poissons} Note: \textit{Sur les surfaces élastiques} gedacht, die im Bulletin, année 1814, S. 47 steht und in deren erstem Teil biegsame, unausdehnbare Flächen betrachtet werden. Die Note ist ein Auszug aus einer Abhandlung, die \textsc{Poisson} am 1. August 1814 gelesen hatte und die in dem zweiten Teil der Mémoires de l’Institut, année 1812, Paris 1816, S. 167 erschienen ist.}.
\par
\textsc{Schumacher} benutzte die erste sich ihm darbietende Gelegenheit und
veranlasste, dass die Kopenhagener Sozietät der Wissenschaften im Jahre 1820
für 1821 die Preisaufgabe stellte, »generaliter superficiem datam in alia superficie
ita exprimere, ut partes minimae imaginis archetypo fiant similes«. Nachdem
keine Abhandlung eingelaufen war, wurde die Aufgabe für 1822 erneuert.
Als \textsc{Schumacher} am 4. Juni 1822 \textsc{Gauss} davon benachrichtigte (Br. G.-Sch. I,
S. 267), antwortete dieser am 10. Juni: »Es tut mir leid, die Wiederholung
Ihrer Preisfrage erst jetzt zu erfahren .\ .\ .\ aber so lange die praktischen
Messungsarbeiten dieses Jahres dauern, kann ich natürlich an eine subtile
theoretische Ausarbeitung gar nicht denken« (Br. G.-Sch. I, S. 270). Am
25. November d. J. fragte er bei seinem Freunde an, bis wann die Preisarbeit
eingesendet werden müsse (Br. G.-Sch. I, S. 293), und nachdem dieser erwiedert
hatte, bis Ende des Jahres, schickte ihm \textsc{Gauss} am 11. Dezember 1825 seine
Ausarbeitung (Br. G.-Sch. I, S. 297). Am 23. Juli 1823 konnte \textsc{Gauss} melden,
\clearpage\SourceLeaf{422}\ArithmeticHead{91}{R}{DIE ALLGEMEINE LEHRE VON DEN KRUMMEN FLÄCHEN.}\setcounter{footnote}{0}
\noindent dass er den Preis erhalten habe (Br. G.-Sch. I, S. 317). Da die Kopenhagener
Gesellschaft sich mit dem Druck der gekrönten Arbeiten nicht befasste, ist
die Abhandlung erst 1825 im dritten und letzten Hefte der von \textsc{Schumacher}
herausgegebenen Astronomischen Abhandlungen erschienen (W. IV, S. 189;
vgl. auch Br. G.-Sch. II, S. 5—7, 17, 22).
\par
Der Begriff der Abbildung steht im Mittelpunkt der \textsc{Gauss}schen Lehre
von den krummen Flächen. »Sie haben ganz Recht«, schreibt \textsc{Gauss} am
11. Dezember 1825 an \textsc{Hansen}, »dass bei allen Kartenprojektionen die Ähnlichkeit
der kleinsten Teile die \textls[100]{wesentliche} Bedingung ist, die man nur in
ganz speziellen Fällen und Bedürfnissen hintansetzen darf. Es wäre wohl
zweckmässig, den Darstellungen, die jener Bedingung Genüge leisten, einen
eigenen Namen zu geben. Inzwischen, allgemein betrachtet, ist sie doch nur
eine Unterabteilung des Generalbegriffs von Darstellung einer Fläche auf einer
andern, die in der Tat \textls[100]{gar nichts weiter} enthält, als dass jedem Punkt der
einen nach irgend einem stetigen \textls[100]{Gesetz} ein Punkt der andern korrespondieren
soll. Es mag wohl etwas Anstrengung kosten, sich zu diesem allgemeinen
Begriff zu erheben; dann aber fühlt man sich auch wirklich auf einem
höhern Standpunkt, wo alles in vergrösserter Klarheit erscheint.\ .\ .\ .\quad Man
kann leicht zeigen, dass, wie allgemein dieser Begriff sei, doch allemal jeder
unendlich kleine Teil (mit Ausnahme der Stellen an singulären Punkten oder
Linien) wahrhaft perspektivisch dargestellt wird, entweder mit völliger Ähnlichkeit,
so wie perspektivische Darstellung auf paralleler Tafel, oder mit halber
Ähnlichkeit, in der in einem Sinn eine Verkürzung statt hat« (Brief im \textsc{Gauss}-Archiv).
\par
Für die Abbildungen, bei denen völlige Ähnlichkeit stattfindet, hat \textsc{Gauss}
im Jahre 1843 das Beiwort \textls[100]{konform} vorgeschlagen (W. IV, S. 262); für den
besonderen Fall der Abbildung des Erdsphäroids auf die Ebene hatte schon
\textsc{Schubert} (1789) von einer projectio conformis gesprochen\footnote{F. Th. \textsc{Schubert}, \textit{De projectione sphaeroidis ellipticae geographica}, Nova acta acad. sc. Petrop., t. 5 ad annum 1787, Petersburg 1789, S. 130.}. Den Satz, dass
eine beliebige stetige Abbildung, von singulären Stellen abgesehen, im Unendlichkleinen
projektiv ist, hat wohl \textsc{Tissot} (1859) zuerst bekannt gemacht\footnote{A. \textsc{Tissot}, \textit{Sur les cartes géographiques}, C. R. t. 49, Paris 1859, S. 673.}.
\par
Die konforme Abbildung hat eine Vorgeschichte. Schon die Griechen
\par\smallskip\noindent\hfill{\small 12\textsuperscript{*}}\par
\clearpage\SourceLeaf{423}\ArithmeticHead{92}{L}{STÄCKEL, GAUSS ALS GEOMETER.}\setcounter{footnote}{0}
\noindent kannten und benutzten die stereographische Projektion der Kugel auf die
Ebene, und \textsc{Gerhard Mercator} (1512—1594) hatte die nach ihm benannte
Abbildung hinzugefügt. \textsc{Lambert} (1772) war dann zu dem allgemeinen Begriff
solcher Abbildungen der Kugel auf eine Ebene gelangt, bei denen die Ähnlichkeit
in den kleinsten Teilen erhalten bleibt, und hatte verschiedene neue
Projektionen dieser Art angegeben, die noch heute bei der Herstellung geographischer
Karten verwendet werden\footnote{J. H. \textsc{Lambert}, \textit{Anmerkungen und Zusätze zur Entwerfung der Land- und Himmelscharten}, erschienen in den \textit{Beyträgen zum Gebrauch der Mathematik}, 3. Teil, Berlin 1772, S. 195—199.}.
\par
\textsc{Lambert} hat in seiner Abhandlung auch die Formeln für die allgemeine
konforme Abbildung einer Kugel auf eine Ebene mitgeteilt, die er \textsc{Lagrange}
verdankte. Dieser geht aus von der bekannten Form für das Quadrat des
Linienelementes der Kugel
\begin{equation}ds^2=dp^2+\cos^2 p\mathbin{.}d\lambda^2,\tag*{(1)}\end{equation}
\noindent die er durch die Substitution
\begin{equation}\mu=\log\operatorname{tang}(45^0+\tfrac12 p)\tag*{(2)}\end{equation}
\noindent auf die Form
\begin{equation}ds^2=\cos^2 p\,(d\lambda^2+d\mu^2)\tag*{(3)}\end{equation}
\noindent bringt. Die Forderung, dass die Kugel konform auf die \(xy\)-Ebene abgebildet
werden soll, führt nunmehr zu der Gleichung
\begin{equation}dx^2+dy^2=\varphi(\lambda,\mu)(d\lambda^2+d\mu^2),\tag*{(4)}\end{equation}
\noindent die, wie die von \textsc{Lagrange} bei Untersuchungen aus der Zahlentheorie häufig
benutzte Identität
\begin{equation}(A^2+B^2)(C^2+D^2)=(AD-BC)^2+(AC+BD)^2\tag*{(5)}\end{equation}
\noindent zeigt, erfüllt ist, wenn man
\begin{equation}dx=nd\lambda-md\mu,\quad dy=md\lambda+nd\mu\tag*{(6)}\end{equation}
\noindent setzt, und es ist daher, wie das \textsc{d’Alembert}sche Verfahren der linearen Verbindungen\footnote{Vgl. P. \textsc{Stäckel}, \textit{Beiträge zur Geschichte der Funktionentheorie im 18. Jahrhundert}, Bibliotheca mathem. (3), 2, 1901, S. 113 und 119.}
erkennen lässt, \(x+\sqrt{-1}\mathbin{.}y\) eine Funktion von \(\lambda+\sqrt{-1}\mathbin{.}\mu\) und
\clearpage\SourceLeaf{424}\ArithmeticHead{93}{R}{DIE ALLGEMEINE LEHRE VON DEN KRUMMEN FLÄCHEN.}\setcounter{footnote}{0}
\noindent gleichzeitig \(x-\sqrt{-1}\mathbin{.}y\) eine Funktion von \(\lambda-\sqrt{-1}\mathbin{.}\mu\). Hieraus ergeben sich
endlich \(x\) und \(y\) als Funktionen von \(\lambda\) und \(\mu\), die der Gleichung (4) genügen.
\par
Bald darauf hat \textsc{Euler} in einer am 4. September 1775 der Petersburger
Akademie vorgelegten, 1778 veröffentlichten Abhandlung denselben Gegenstand
behandelt\footnote{L. \textsc{Euler}, \textit{De repraesentatione superficiei sphaericae super plano}, Acta acad. sc. Petrop. t. 1 pro anno 1777; I, 1778, S. 107.}. Er vermeidet den Kunstgriff, die Identität (5) heranzuziehen,
und gewinnt die Gleichungen (6) oder doch mit ihnen gleichbedeutende Gleichungen
unmittelbar aus der Forderung der Ähnlichkeit in den kleinsten
Teilen. Damit ist zugleich bewiesen, dass das Bestehen der Gleichungen (6)
nicht nur, wie bei \textsc{Lagrange}, hinreichend, sondern auch notwendig ist. Während
ferner \textsc{Lambert} aus den Formeln von \textsc{Lagrange} keinen Nutzen gezogen
hatte, gelingt es \textsc{Euler}, mit ihrer Hilfe besondere Lösungen der Aufgabe
herzuleiten.
\par
Ein Blick auf die vorstehenden Formeln lässt erkennen, dass das Verfahren
von \textsc{Lagrange} und \textsc{Euler} sich ohne weiteres auf den allgemeineren
Fall übertragen lässt, wo das Quadrat des Linienelementes der krummen Fläche,
die konform auf die Ebene abgebildet werden soll, auf die Gestalt
\begin{equation}ds^2=\psi(\lambda,\mu)(d\lambda^2+d\mu^2)\tag*{(7)}\end{equation}
\noindent gebracht werden kann. Für die Drehflächen, bei denen vermöge der Meridiane
und Parallelkreise
\begin{equation}ds^2=dp^2+G(p)d\lambda^2\tag*{(8)}\end{equation}
\noindent ist, gelingt das sofort durch die Substitution
\begin{equation}\mu=\int\frac{dp}{\sqrt{G(p)}}\mathpunct{\raisebox{.5ex}{.}}\tag*{(9)}\end{equation}
\noindent Auf diese Weise ist \textsc{Lagrange} 1781 zu den allgemeinen Formeln für die
konforme Abbildung einer Drehfläche auf eine Ebene gelangt; er hat davon
schöne Anwendungen gemacht\footnote{J. L. \textsc{Lagrange}, \textit{Sur la construction des cartes géographiques}, Nouv. Mém. de l’Acad., année 1779, Berlin 1781, S. 161, 186; Oeuvres, t. 4, S. 635.}.
\par
Der Fortschritt, den \textsc{Gauss} in der Preisschrift vom Jahre 1822 gemacht
hat, liegt darin, dass er zeigte, wie man bei einer \textls[100]{beliebigen} (reellen)
krummen Fläche, bei der das Quadrat des Linienelements in der allgemeinen
\clearpage\SourceLeaf{425a}\ArithmeticHead{94}{L}{STÄCKEL, GAUSS ALS GEOMETER.}\setcounter{footnote}{0}
\noindent Form
\begin{equation}ds^2=E\,dp^2+2F\,dp\,dq+G\,dq^2\tag*{(10)}\end{equation}
\noindent gegeben ist, die besondere Form
\begin{equation}ds^2=\psi(\lambda,\mu)(d\lambda^2+d\mu^2)\tag*{(11)}\end{equation}
\noindent herstellen kann. Dies geschieht, indem die Gleichung
\begin{equation}ds^2=0\tag*{(12)}\end{equation}
\noindent integriert wird, oder was auf dasselbe hinauskommt, indem für zwei konjugiert
komplexe lineare Differentialformen, deren Produkt \(ds^2\) ist, die ebenfalls konjugiert
komplexen \textsc{Euler}schen Multiplikatoren ermittelt werden. Damit aber
erhält man zugleich die \textls[100]{allgemeine} konforme Abbildung der gegebenen
krummen Fläche auf die Ebene, sodass es des Verfahrens von \textsc{Lagrange} gar
nicht mehr bedarf, und während bei \textsc{Lagrange} das Imaginäre nur formal als
Mittel zur Integration der Gleichungen (6) auftrat, ist jetzt durch die Gleichung
\(ds^2=0\) der wahre Grund für das Auftreten von Funktionen komplexer Grössen
aufgedeckt.
\par
Wer sich der hier dargelegten Auffassung anschliesst, wird dem Urteil
\textsc{Jacobis} nicht beipflichten können, dass »der \textsc{Lagrange}schen Arbeit nur wenig
hinzuzusetzen war«\footnote{C. G. \textsc{Jacobi}, \textit{Vorlesungen über Dynamik}, gehalten im W.-S. 1842/43, 2. Ausgabe, Berlin 1884, S. 215; vgl. auch Br. G.-Sch., III, S. 173, der beim Abdruck unterdrückte Name ist v. \textsc{Littrow}.}. \textsc{Jacobi} hat auch beanstandet, dass \textsc{Gauss} diese Arbeit
nicht erwähnt habe; allein \textsc{Gauss} ist überhaupt nicht auf die Geschichte der
konformen Abbildung eingegangen, ebenso wie sich auch \textsc{Euler} aller Anführungen
enthalten hatte.
\clearpage\endgroup
